\begin{titlepage}
    \centering
    \vspace*{1.8cm}
    {\zihao{2}\bfseries SPH 平板泊肃叶流算例技术报告\par}
    \vspace{1.2cm}
    {\zihao{4}项目：\code{SPH-Poiseuille-Flow}\par}
    \vspace{0.5cm}
    {\zihao{-4}交付内容：LaTeX 源文件、编译 PDF、公式与代码映射说明\par}
    \vspace{3.5cm}
    \begin{tabular}{rl}
        报告对象： & 二维弱可压 SPH 平板泊肃叶流单层 shell 无滑移算例 \\
        核心脚本： & \code{SPH\_Poiseuille.m} \\
        关键 MEX： & \code{sph\_neighbor\_search\_mex}、\code{sph\_physics\_shell\_mex} \\
        生成日期： & 2026-03-19 \\
    \end{tabular}
    \vfill
    {\zihao{-4}本报告面向项目维护、数值方法复核与后续算例扩展。\par}
\end{titlepage}

\pagenumbering{Roman}

\chapter*{摘要}
\addcontentsline{toc}{chapter}{摘要}
本报告围绕仓库 \code{SPH-Poiseuille-Flow} 的当前实现展开，目标是把算例中的物理模型、SPH 离散、单层 CFL 时间推进、边界处理、MEX 加速结构和验证流程整理成一份可以直接审阅与复现的技术文档。报告以源码为唯一实现依据，以 \code{README.md}、\code{config.ini}、\code{build\_shell\_wall\_particles.m} 以及现有结果图为辅助材料，将连续方程、弱可压状态方程、核函数、密度重初始化、核梯度修正、粘性项、压力项、transport correction、两阶段积分和单层 shell 壁面逐一映射到具体文件与代码段。

该算例的计算对象是二维平板通道中的重力驱动层流。主程序 \code{SPH\_Poiseuille.m} 现已整理为“单主文件壳层 + local functions”的形式，负责配置读取、粒子初始化、全局调度、restart 管理和后处理，并在运行时把数据组织为 \code{cfg}、\code{state}、\code{geom}、\code{neighbor} 和 \code{monitor} 五类结构；\code{build\_shell\_wall\_particles.m} 负责生成上下单层 shell 壁面粒子的几何元数据；\code{sph\_neighbor\_search\_mex.c} 负责周期性邻居搜索；\code{mex/sph\_physics\_mex.c} 则在运行时被编译为 \code{sph\_physics\_shell\_mex}，同时提供底层物理 mode 与两个高层门面接口 \code{advance\_shell\_step}、\code{wall\_shear\_monitor}。整体结构依然保持 MATLAB 与 C/MEX 分工协作：MATLAB 保持可读性和流程控制，MEX 将粒子对上的高频循环下沉到 C 层。

最近一次提交已经把旧的多方案壁面逻辑收敛为“单层 shell 壁面 + 唯一 no-slip”实现，同时删除了 \code{bc\_mode}、PI 控制、Neumann/ghost 分支和独立的 \code{resolution\_study.m}。因此，本报告同步修正文档口径，明确以下事实：代码实际使用二维 cubic spline 核；\code{c\_f} 在实现中直接作为人工声速参与状态方程和统一 CFL；transport correction 的实现系数为写死的 $0.2h^2$；当前验证链条由主结果图、中截面剖面演化和运行期壁面切应力日志共同构成。通过这些澄清，读者可以安全地理解“仓库当前版本真正做了什么”。

\paragraph{关键词}
泊肃叶流；弱可压 SPH；单层 shell 壁面；无滑移边界；单层 CFL 时间步进；MATLAB MEX

\chapter*{符号与缩写}
\addcontentsline{toc}{chapter}{符号与缩写}
\begin{longtable}{p{0.18\textwidth} p{0.72\textwidth}}
\toprule
符号或缩写 & 含义 \\
\midrule
\endhead
$DL$ & 通道长度，沿 $x$ 方向采用周期边界 \\
$DH$ 或 $H$ & 通道高度，沿 $y$ 方向由上下壁面约束 \\
$dp$ & 初始粒子间距，决定空间分辨率 \\
$h$ & 光滑长度，代码中取 $h=1.3dp$ \\
$\rho_0$ & 参考密度 \\
$\mu$ & 动力粘度，$\nu=\mu/\rho_0$ 为运动粘度 \\
$U_{bulk}$ & 目标截面平均速度 \\
$U_{max}$ & 理论中心线最大速度，二维平板泊肃叶流中 $U_{max}=1.5U_{bulk}$ \\
$g$ & 等效体积力或压力梯度驱动力 \\
$p_0$ & 参考压力，代码中取 $p_0=\rho_0 c^2$ \\
$c_f$ & 当前实现中直接作为人工声速使用的参数名 \\
$W_{ij}$ & 粒子对 $(i,j)$ 的核函数值 \\
$\partial W_{ij}/\partial r$ & 核函数对径向距离的导数 \\
$V_i$ & 粒子体积，离散中取 $V_i=m_i/\rho_i$ \\
$B_i$ & 核梯度修正矩阵或其扁平化表示 \\
$dt$ & 统一 CFL 时间步 \\
$\tau_w$ & 解析壁面切应力，用于运行期监测与验证 \\
shell 壁面 & 中面位于通道外侧、显式携带厚度的单层壁面粒子 \\
MEX & MATLAB Executable，MATLAB 调用的 C 扩展模块 \\
\bottomrule
\end{longtable}

\tableofcontents
\clearpage
